% ===== 7. Критерії і ваги =====
\begin{frame}{Критерії і ваги}
  \centering
  \begin{tabular}{@{}l c l@{}}
    \toprule
    \textbf{Критерій} & \textbf{Вага} & \textbf{Що означає} \\
    \midrule
    Актуальність проблеми        & 0,25 & Наскільки проблема важлива для аудиторії \\
    Реалізованість за семестр    & 0,25 & Чи зробимо робочий продукт \\
    Конкурентна перевага         & 0,2  & Чим ми кращі за наявні рішення \\
    Інтерес і досвід команди     & 0,15 & Чи хочемо і чи маємо потрібні знання \\
    Ризики і зовнішні залежності & 0,15 & Менше ризиків — вища оцінка \\
    \midrule
    \textbf{Разом}               & \textbf{1,0} & \\
    \bottomrule
  \end{tabular}

  \note{[Доповідач 3] Щоб порівняти теми чесно, ми домовилися про п'ять критеріїв. Найбільшу вагу, по 0,25, мають актуальність проблеми і реалізованість за семестр. Тобто продукт має бути потрібним, і ми маємо встигнути його зробити. Конкурентна перевага важить 0,2: чим ми кращі за те, що вже є. Інтерес і досвід команди та ризики — по 0,15. Для ризиків шкала обернена: що менше ризиків і зовнішніх залежностей, то вища оцінка. Разом ваги дають одиницю, тож підсумок — це зважена оцінка на тій самій шкалі від одного до п'яти.}
\end{frame}

% ===== 8. Порівняння =====
\begin{frame}{Порівняння тем}
  \centering\renewcommand{\arraystretch}{1.1}
  \begin{tabular}{@{}l c >{\columncolor{highlight}}c c@{}}
    \toprule
    \textbf{Критерій} & \textbf{Вага} & \textbf{\begin{tabular}{@{}c@{}}Цифровий\\тренер\end{tabular}} & \textbf{\begin{tabular}{@{}c@{}}Розумний\\кошик\end{tabular}} \\
    \midrule
    Актуальність проблеми        & 0,25 & 4 & 4 \\
    Реалізованість за семестр    & 0,25 & 4 & 4 \\
    Конкурентна перевага         & 0,2  & 3 & 2 \\
    Інтерес і досвід команди     & 0,15 & ? & ? \\
    Ризики і зовнішні залежності & 0,15 & 3 & 2 \\
    \midrule
    \textbf{Зважена сума без «інтересу»} (макс. 4,25) & 0,85 & \textbf{3,05} & \textbf{2,70} \\
    \bottomrule
  \end{tabular}

  \vspace{0.3em}
  \raggedright\small
  Підсумок = зважена сума + 0,15 × середня оцінка інтересу команди.\\
  Кошик випередить тренера, лише якщо його інтерес вищий хоча б на 3 бали.

  \note{[Доповідач 3] Ось самі оцінки. За актуальністю і реалізованістю теми рівні — обидві по 4. Різниця з'являється в конкурентній перевазі: 3 у тренера проти 2 у кошика. І в ризиках теж 3 проти 2. Оцінки за інтерес команди ми ще не зібрали, тому там знаки питання, і зважену суму рахуємо без цього критерію: максимум тут 4,25. Тренер має 3,05, кошик — 2,70, різниця 0,35 бала. Коли кожен з нас поставить свою оцінку, ми додамо 0,15, помножене на середнє. Кошик виграв би, лише якби ми оцінили його інтерес хоча б на 3 бали вище.}
\end{frame}
